%; whizzy paragraph -pdf xpdf -latex ./whizzypdfptex.sh
%; whizzy-paragraph "^\\\\begin{frame}"
% latex beamer presentation.
% platex, latex-beamer でコンパイルすることを想定。 

%     Tokyo Debian Meeting resources
%     Copyright (C) 2009 Junichi Uekawa
%     Copyright (C) 2009 Nobuhiro Iwamatsu

%     This program is free software; you can redistribute it and/or modify
%     it under the terms of the GNU General Public License as published by
%     the Free Software Foundation; either version 2 of the License, or
%     (at your option) any later version.

%     This program is distributed in the hope that it will be useful,
%     but WITHOUT ANY WARRANTY; without even the implied warreanty of
%     MERCHANTABILITY or FITNESS FOR A PARTICULAR PURPOSE.  See the
%     GNU General Public License for more details.

%     You should have received a copy of the GNU General Public License
%     along with this program; if not, write to the Free Software
%     Foundation, Inc., 51 Franklin St, Fifth Floor, Boston, MA  02110-1301 USA

\documentclass[cjk,dvipdfmx,12pt]{beamer}
\usetheme{Tokyo}
\usepackage{monthlypresentation}

%  preview (shell-command (concat "evince " (replace-regexp-in-string
%  "tex$" "pdf"(buffer-file-name)) "&")) 
%  presentation (shell-command (concat "xpdf -fullscreen " (replace-regexp-in-string "tex$" "pdf"(buffer-file-name)) "&"))
%  presentation (shell-command (concat "evince " (replace-regexp-in-string "tex$" "pdf"(buffer-file-name)) "&"))

%http://www.naney.org/diki/dk/hyperref.html
%日本語EUC系環境の時
\AtBeginDvi{\special{pdf:tounicode EUC-UCS2}}
%シフトJIS系環境の時
%\AtBeginDvi{\special{pdf:tounicode 90ms-RKSJ-UCS2}}

\title{Debian / Ubuntu ユーザーミートアップ in 札幌 2017.07 \\　\\lxcについて}
\subtitle{}
\author{Norimitsu Sugimoto (杉本 典充) \\dictoss@live.jp}
\date{2017-07-14}
\logo{\includegraphics[width=8cm]{image200607/openlogo-light.eps}}

\begin{document}

\frame{\titlepage{}}

%\section{}

\begin{frame}{アジェンダ}
  \begin{itemize}
  \item 自己紹介
  \item 仮想化技術について
  \item lxcとは
  \item lxcのインストール
  \item lxcのコマンド解説
  \item lxcを実用する
  \item おわりに
  \item 参考資料
  \end{itemize}
\end{frame}


\begin{frame}{自己紹介}
  \begin{itemize}
  \item Norimitsu Sugimoto (杉本 典充)
  \item dictoss@live.jp
  \item Twitter: @dictoss
  \item Debian-3.1、FreeBSD-6.2の頃から使っています
  \item Debian GNU/kFreeBSDが気になっておりウォッチ中
  \item 仕事はソフトウェア開発者をやってます
  \end{itemize}
\end{frame}


\emtext{仮想化技術について}

\begin{frame}[containsverbatim]{仮想化技術の分類}
  \begin{itemize}
  \item コンテナ型仮想化
  \item 準仮想化型
  \item 完全仮想化型 (エミュレーション型)
  \item 完全仮想化型 (ハイパーバイザ型)
  \end{itemize}
\end{frame}

\begin{frame}[containsverbatim]{仮想化技術のメリット・デメリット}
  \begin{itemize}
  \item 準仮想化型、完全仮想化型
    \begin{itemize}
    \item 物理マシンをエミュレートした仮想マシンとして動作する
    \item 仮想マシン上でもカーネルを動作させる
    \item 物理マシンで動かしていたプログラムはほぼそのまま動く
    \item CPU、メモリ、ディスクを多く消費する
    \end{itemize}
  \item コンテナ型仮想化
    \begin{itemize}
    \item ゲスト環境の動作にカーネルは不要(ホスト環境のカーネル上で動作)
    \item ゲスト環境は、ホスト環境から見るとプロセスとして扱われる
    \item ゲスト環境が利用できるリソースに制約がつく場合がある
    \end{itemize}
  \end{itemize}
\end{frame}

\begin{frame}[containsverbatim]{chroot}
  \begin{itemize}
  \item chrootシステムコールとchrootコマンド
  \item 1982年にビル・ジョイが開発したとされている
  \item \# chroot rootfsdir でコンテナ環境に入れることができる
  \end{itemize}
\end{frame}


\emtext{lxc}

\begin{frame}[containsverbatim]{lxcとは}
  \begin{itemize}
  \item LinuX Containersのことで、省略してlxcと読んでいる
  \item あるディレクトリ配下に実行ファイル、ライブラリ、設定ファイルを適切に配置したrootfsを準備する
  \item rootfsをchroot環境で起動し、仮想マシンのように動かすことができる
  \item Debian 9 Stretch では lxc-2.0.5 を採用
  \end{itemize}
\end{frame}


\emtext{lxcのインストール}

% https://wiki.debian.org/LXC/LibVirtDefaultNetwork
\begin{frame}[containsverbatim]{インストールの流れ}
  \begin{itemize}
  \item ブリッジネットワークとlibvirtdの準備
  \item lxcのインストール
  \end{itemize}
\end{frame}


\begin{frame}[containsverbatim]{ブリッジネットワークとlibvirtdの準備}

パッケージのインストール
\begin{commandline}
  # apt-get install libvirt-clients \
  libvirt-daemon-system ebtables dnsmasq
\end{commandline}

仮想ネットワークの設定
\begin{commandline}
  # virsh net-autostart default
  # virsh net-start default
\end{commandline}

\end{frame}


\begin{frame}[containsverbatim]{ブリッジネットワークとlibvirtdの準備}
確認
\begin{commandline}
  $ sudo virsh net-info default
  Name:           default
  UUID:           78564864-f237-4059-a12a-3ec04369a27b
  Active:         yes
  Persistent:     yes
  Autostart:      yes
  Bridge:         virbr0
  
  $ ip a show virbr0
  →192.168.122.1/24 が付与されている
\end{commandline}

\end{frame}


\begin{frame}[containsverbatim]{lxcのインストール}
  コンテナ内のリソース制約を処理するcgroupの確認
  \begin{commandline}
    # mount | grep cgroup
    →stretchは標準でmountされている
  \end{commandline}
  パッケージのインストール
  \begin{commandline}
    # apt-get install lxc libvirt0 libpam-cgroup \
    libpam-cgfs
  \end{commandline}
  環境を確認
  \begin{commandline}
  $ ls /usr/bin | grep lxc
  # lxc-checkconfig
  \end{commandline}
\end{frame}

\begin{frame}[containsverbatim]{lxcのインストール}
  \begin{itemize}
  \item lxc-create
  \item lxc-destroy
  \item lxc-start
  \item lxc-stop
  \item lxc-console
  \item lxc-attach
  \end{itemize}
\end{frame}


\emtext{lxcのコマンド解説}

\begin{frame}[containsverbatim]{lxc-create(1)}
  \begin{commandline}
    # lxc-create -n demo1 -t debian -- \
    --release=stretch --arch=amd64 \
    --mirror=http://ftp.jp.debian.org/debian
  \end{commandline}   
  \begin{itemize}
  \item 実行するとテンプレートがdebianの場合はdebootstrapを実行してrootfsをダウンロードする
  \item lxcのゲスト環境のディレクトリは、/var/lib/lxc/''コンテナ名'' 。中身は以下。
    \begin{itemize}
    \item config (設定ファイル)
    \item rootfs (コンテナの中身)
    \end{itemize}
  \end{itemize}
\end{frame}

\begin{frame}[containsverbatim]{lxc-create(2)}
  \begin{itemize}
  \item configを修正して、ネットワークの設定を行う
  \item lxc.network.type = veth
  \item lxc.network.flags = up
  \item lxc.network.link = virbr0
  \item lxc.network.name = eth0
  \item lxc.network.ipv4 = 192.168.122.60/24
  \item lxc.network.ipv4.gateway = 192.168.122.1
  \end{itemize}
\end{frame}

\begin{frame}[containsverbatim]{lxc-destroy}
  \begin{commandline}
  # lxc-destroy -n demo1
  \end{commandline}
  \begin{itemize}
  \item コンテナを削除します
  \end{itemize}
\end{frame}

\begin{frame}[containsverbatim]{lxc-start}
  \begin{commandline}
  # lxc-start -n demo1
  \end{commandline}
  \begin{itemize}
  \item 実行して何もエラーが表示されなければ、バックグラウンドでlxcコンテナが動き出します
  \item 起動したコンテナへの接続は、後述するlxc-consoleまたはlxc-attachで行います
  \item コンテナへのログインはsshでもログインできますが、ユーザを作成する必要があります
  \end{itemize}
\end{frame}

\begin{frame}[containsverbatim]{lxc-stop}
  \begin{commandline}
  # lxc-stop -n demo1
  \end{commandline}
  \begin{itemize}
  \item コンテナ環境を終了するよう指示を出します
  \item コンテナ環境の終了とは、コンテナ内のinitプログラムを終了することをいいます
  \item コンテナ環境でshutdown命令は実行できません
  \end{itemize}
\end{frame}

\begin{frame}[containsverbatim]{lxc-console}
  \begin{commandline}
  # lxc-console -n demo1
  \end{commandline}
  \begin{itemize}
  \item lxcのゲスト環境のコンソールに接続します
  \item コンソールを抜ける場合は、「Ctrl+a q」の順に入力してください
  \end{itemize}
\end{frame}

\begin{frame}[containsverbatim]{lxc-attach}
  \begin{commandline}
  # lxc-attach -n demo1 {command}
  \end{commandline}
  \begin{itemize}
  \item lxcのゲスト環境でコマンドを実行します
  \item コマンドを指定しない場合は、コンテナ内のユーザのデフォルトシェルが実行されます
  \item lxc-consoleでログインすることに比べ、いきなりコンテナ内でシェルを実行できるため、lxc-attachの方がコンテナ内の整備がしやすいです
  \end{itemize}
\end{frame}


\emtext{lxcを実用する}

\begin{frame}[containsverbatim]{コンテナ環境のセットアップの流れ}
  \begin{itemize}
  \item lxc-create を実行してコンテナを生成する
  \item lxcのゲスト環境のconfigを書き換えてネットワークを設定する
  \item lxc-startしてコンテナを起動する
  \item lxc-attachでゲスト環境に入る
  \begin{commandline}
  # passwd
  # adduser username
  # apt-get install sudo vim-tiny
  # visudo
  \end{commandline}
  \item sshログインしてお好みに設定する
  \end{itemize}
\end{frame}

\begin{frame}[containsverbatim]{何にlxcを使うか}
  \begin{itemize}
  \item 一時的な検証で、ホスト環境にいろいろインストールしたくない場合
  \item アプリケーションのクリーンビルドやクリーンインストールをテストする場合
  \item ホスト環境はsystemd、ゲスト環境はsysvinitと使い分ける場合
  \item python2系とpython3系のwsgiアプリを1つのホストで動かしたい場合
  \item ホスト環境と異なるCPUアーキテクチャのエミュレーション環境がほしい場合
    \begin{itemize}
    \item apt-get install qemu qemu-user-static binfmt-support
    \item その後、lxc-createを実行してください
    \item 詳しくは CrossDebootstrap を調べてみてください
    \end{itemize}
  \end{itemize}
\end{frame}


\begin{frame}[containsverbatim]{おわりに}
  \begin{itemize}
  \item Debian上でlxcを試してみました
  \item 発展系であるLXDやdockerへつなげていきましょう
  \item コンテナは便利ですので試してみてください
  \end{itemize}
\end{frame}


\begin{frame}{参考情報}
  \begin{itemize}
  \item 「LXC」 \url{https://linuxcontainers.org/}
  \item 「LXC - Debian Wiki」 \url{https://wiki.debian.org/LXC}
  \item 「LXCLibVirtDefaultNetwork」 \url{https://wiki.debian.org/LXC/LibVirtDefaultNetwork}
  \end{itemize}
\end{frame}

\end{document}


;;; Local Variables: ***
;;; outline-regexp: "\\([ 	]*\\\\\\(documentstyle\\|documentclass\\|emtext\\|section\\|begin{frame}\\)\\*?[ 	]*[[{]\\|[]+\\)" ***
;;; End: ***
